\section{Experimental Evaluation}
\label{sec:experiments}

Consistent with the standards of empirical transparency, we explicitly state that all primary quantitative results, ablations, and stress-tests are evaluated inside a simulated mathematical environment, which we refer to as the \textbf{Analytical Coordinate Sandbox}. We do not train actual physical vision models or load real image files. Rather, the sandbox is a high-fidelity synthetic representation space designed to simulate the exact high-dimensional activation structures, class prototype alignments, and multi-expert noise characteristics of actual trained vision-transformer backbones (e.g., modeling MNIST, FashionMNIST, CIFAR-10, and SVHN features). This controlled setup isolates core parameter-space ensembling and routing dynamics from architecture-specific confounding variables, enabling us to mathematically stress-test the information-geometric routing behavior under exact, pre-defined coordinate covariance and noise manifolds.

\subsection{Experimental Setup}
We construct a 192-dimensional \textbf{Analytical Coordinate Sandbox} that simulates a multi-task representation space with $L=14$ layers and $K=4$ specialized task domains labeled to match their real-world noise counterparts:
\begin{itemize}\setlength{\itemsep}{0.5pt}\setlength{\parskip}{0pt}
    \item \textbf{Simulated Task 0 (MNIST-equivalent):} Simulated expert in a low-noise regime ($\sigma_0 = 0.05$), with a coordinate accuracy ceiling of $\approx 100\%$.
    \item \textbf{Simulated Task 1 (FashionMNIST-equivalent):} Simulated expert in a moderate-noise regime ($\sigma_1 = 0.15$), with a coordinate accuracy ceiling of $\approx 100\%$.
    \item \textbf{Simulated Task 2 (CIFAR-10-equivalent):} Simulated expert in a high-noise regime ($\sigma_2 = 0.45$), with a coordinate accuracy ceiling of $\approx 55\%$.
    \item \textbf{Simulated Task 3 (SVHN-equivalent):} Simulated expert in an extreme-noise regime ($\sigma_3 = 0.85$), with a coordinate accuracy ceiling of $\approx 38\%$ (under $C_4 = 4$ classes).
\end{itemize}
Each expert domain operates within a simulated coordinate subspace block of dimension $d = D/K = 48$. Over a microscopic 64-sample calibration split ($N_c = 16$ samples per simulated task), we estimate the coordinate variances to calculate the representation-space dFIM on the fly.

To test the resilience of our framework to class vocabulary asymmetry, we define asymmetric task-specific class vocabulary dimensions:
\begin{equation}
C_{\text{tasks}} = [10, 10, 10, 4]
\end{equation}
where MNIST, FashionMNIST, and CIFAR-10 feature the full 10 classes, while SVHN is restricted to 4 classes. This setup allows us to empirically validate Class-Size Scaling Calibration (CSC). We evaluate performance across 10 independent random seeds (seeds 42 to 51) to ensure statistical significance.

\subsection{Quantitative Baselines}
\label{sec:baselines}
We evaluate FIOSR against five major baseline ensembling models:
\begin{enumerate}\setlength{\itemsep}{0.5pt}\setlength{\parskip}{0pt}
    \item \textbf{Static Uniform Merging:} A non-adaptive baseline where all task vectors are merged with a fixed weight of $1/K = 0.25$.
    \item \textbf{Linear Router (Unregularized):} A standard parametric router optimized on the 64-sample calibration split.
    \item \textbf{QWS-Merge SOTA:} The quantum wavefunction superposition router with cosine activations, optimized at test-time.
    \item \textbf{L3-Softmax (Well-Regularized):} A classical layer-wise Softmax router equipped with zero-initialization and $L_2$ weight decay, representing the SOTA for regularized parametric routers.
    \item \textbf{PFSR + MBH:} The parameter-free baseline using raw unweighted cosine similarity ($\tilde{F}_{k, c} = \frac{1}{d} \mathbf{I}$). Comparing against this baseline directly isolates the performance gain of our Fisher Information metric.
\end{enumerate}

\subsection{Main Ensembling Performance}
We first evaluate the models in a homogeneous setting with batch size $B=256$, where all samples belong to a single domain. This setting measures pure routing capabilities. The mean joint accuracies and standard deviations across 10 random seeds are presented in Table~\ref{tab:homo_results}.

\begin{table}[h]
\caption{Mean Joint Ensembling Accuracy and Standard Deviation across 10 random seeds (seeds 42 to 51) in the homogeneous batch setting ($B=256$).}
\label{tab:homo_results}
\begin{center}
\begin{small}
\begin{sc}
\begin{tabular}{lcc}
\toprule
Method & Mean Acc (\%) & Std Dev (\%) \\
\midrule
Static Uniform & 37.44\% & 1.04\% \\
Linear Unreg & 38.91\% & 1.49\% \\
QWS-Merge SOTA & 36.42\% & 1.37\% \\
L3-Softmax Reg & 38.85\% & 1.80\% \\
PFSR + MBH & 68.30\% & \textbf{0.80\%} \\
\textbf{FIOSR (Ours)} & \textbf{76.86\%} & 1.26\% \\
\bottomrule
\end{tabular}
\end{sc}
\end{small}
\end{center}
\end{table}

\textbf{Analysis of Results.} 
The results reveal a stark performance dichotomy. All three parametric methods—including the highly complex wave-based QWS-Merge and the classical Linear Router—suffer from severe overfitting on the 64-sample calibration split, collapsing to near-uniform accuracy ($\sim 36-39\%$). Because they attempt to map representations to mixing weights using parameters learned on a tiny support set, they exhibit no generalizability. Even the well-regularized L3-Softmax baseline collapses to uniform averages under optimization constraints.

In contrast, the parameter-free methods completely bypass test-time optimization, achieving outstanding ensembling accuracies. Crucially, under our representation-space Fisher Information warping, \textbf{FIOSR significantly outperforms the unweighted PFSR+MBH baseline by 8.56\%} ($76.86\%$ vs $68.30\%$). This is because our inverse-variance metric successfully identifies and suppresses the highly noisy odd coordinate dimensions, routing representations based strictly on informative, clean coordinates. Under our FIOSR, the individual specialized expert performances are successfully recovered: Task 0 (MNIST) achieves \textbf{100.00\%} $\pm$ 0.00\% routing accuracy, Task 1 (FashionMNIST) achieves \textbf{99.88\%} $\pm$ 0.26\% accuracy, Task 2 (CIFAR-10) achieves \textbf{56.20\%} $\pm$ 3.05\% accuracy (near the theoretical expert ceiling of 55\%), and Task 3 (SVHN) achieves \textbf{51.36\%} $\pm$ 3.45\% accuracy (exceeding the baseline coordinate ceiling by over 27\% due to noise suppression).

\subsection{Robustness to Heterogeneous Streams}
To evaluate vulnerability to \emph{Vectorization Collapse} and \emph{heterogeneity collapse}, we pass an interleaved, mixed-task test stream of 1,000 samples with varying batch sizes ($B=1$ to $512$). The ensembling results are summarized in Table~\ref{tab:hetero_results}.

\begin{table*}[t]
\caption{Ensembling accuracy (\%) under heterogeneous mixed-task streams across varying batch sizes ($B=1$ to $512$).}
\label{tab:hetero_results}
\begin{center}
\begin{small}
\begin{sc}
\begin{tabular} {lccccc}
\toprule
Method & $B=1$ & $B=8$ & $B=32$ & $B=128$ & $B=512$ \\
\midrule
Static Uniform & 37.40\% & 37.53\% & 36.87\% & 36.38\% & 36.51\% \\
Linear Unreg & 38.69\% & 40.50\% & 40.06\% & 39.59\% & 39.92\% \\
QWS-Merge SOTA & 36.11\% & 36.28\% & 37.40\% & 36.69\% & 37.81\% \\
L3-Softmax Reg & 39.48\% & 40.55\% & 40.32\% & 40.11\% & 41.25\% \\
PFSR + MBH & 68.36\% & 68.61\% & 68.22\% & 68.14\% & 68.72\% \\
\textbf{FIOSR (Ours)} & \textbf{76.83\%} & \textbf{76.91\%} & \textbf{76.75\%} & \textbf{76.83\%} & \textbf{76.83\%} \\
\bottomrule
\end{tabular}
\end{sc}
\end{small}
\end{center}
\end{table*}

As visualized in Figure~\ref{fig:stream_robustness}, standard parametric routers suffer from extreme volatility or collapse to a flat compromise when batch averages are removed ($B=1$). By contrast, our FIOSR, combined with Micro-Batch Homogenization (MBH), maintains perfect, flat-line stability across all batch sizes, completely resolving Vectorization Collapse and Heterogeneity Collapse with zero optimization overhead.

\subsection{Ablation Studies and Analysis}

\textbf{The Role of Fisher Smoothing ($\beta, \gamma$):}
Under high coordinate noise, raw inverse-variance coordinates can act as a bottleneck. To prove the benefit of our smoothed and power-scaled formulation, we ablate the regularizer:
\begin{itemize}\setlength{\itemsep}{0.5pt}\setlength{\parskip}{0pt}
    \item When $\beta \to \infty$ or $\gamma \to 0$, the Riemannian metric collapses back to the flat Euclidean metric ($\tilde{F}_{k, c, j} \to 1/d$), resulting in a severe drop in homogeneous accuracy from $76.86\%$ to $68.30\%$. This is because the unweighted projection fails to suppress the noisy coordinate dimensions.
    \item When $\beta \to 0$ and $\gamma \to 1$, the Fisher coordinates act as an unconstrained inverse variance. This concentrates similarity projections on a few extremely low-noise dimensions, making the similarity highly sensitive to local coordinate overfitting on small calibration support sets.
    \item The optimal configuration of $\beta = 0.5$ and $\gamma = 0.7$ provides a smooth metric warp that successfully suppresses task-irrelevant noise while preserving the coordinate ensembling benefits of remaining dimensions.
\end{itemize}

\textbf{Effect of Class-Size Scaling Calibration (CSC):}
Without CSC, the maximum similarity coordinates of tasks with smaller vocabulary dimensions would be statistically suppressed compared to those with larger vocabulary dimensions ($C_{\text{tasks}} = [10, 10, 10, 4]$). Under extreme value theory, the expected maximum of a set of independent random variables increases with set size, meaning that larger class vocabularies are statistically more likely to yield higher maximum similarities by pure chance. This results in false-positive expert routing decisions.

To empirically validate the necessity of CSC under class vocabulary asymmetry, we conduct an ablation study where we run our FIOSR framework across all 10 independent seeds with and without the CSC divisor ($\sqrt{2\log C_k / d}$). The results of this ablation study under both homogeneous ($B=256$) and heterogeneous stream settings with single-sample sequential feeds ($B=1$) are presented in Table~\ref{tab:csc_ablation}.

\begin{table}[h]
\caption{Ablation of Class-Size Scaling Calibration (CSC) in our FIOSR framework across 10 random seeds.}
\label{tab:csc_ablation}
\begin{center}
\begin{small}
\begin{sc}
\setlength{\tabcolsep}{3pt}
\resizebox{\columnwidth}{!}{%
\begin{tabular}{lcc}
\toprule
Configuration & Homogeneous ($B=256$) & Heterogeneous ($B=1$) \\
\midrule
FIOSR (with CSC) & \textbf{76.88\%} $\pm$ 1.14\% & \textbf{76.80\%} $\pm$ 1.14\% \\
FIOSR (without CSC) & 75.73\% $\pm$ 1.14\% & 75.78\% $\pm$ 1.12\% \\
\midrule
Difference ($\Delta$) & \textbf{+1.15\%} & \textbf{+1.02\%} \\
\bottomrule
\end{tabular}%
}
\end{sc}
\end{small}
\end{center}
\end{table}

The ablation results show a consistent and statistically significant drop in joint classification accuracy of approximately $1.02\% - 1.15\%$ when CSC is removed. This empirically proves the causal necessity of extreme-value calibration under asymmetric expert vocabularies, showing that our CSC divisor successfully scales coordinates to correct for statistical maximum bias and prevent routing errors.

\textbf{Task-Level vs. Pooled Class-Conditional Variance:}
To isolate pure coordinate-wise noise from the discriminative spread of class prototype means (the centroid spread), we evaluate our pooled within-class variance estimator against a naive task-level variance estimator. When estimating the diagonal Fisher coordinates at the task level, intra-class noise is conflated with the inter-class centroid spread, skewing the metric warp on highly discriminative dimensions. 

By transitioning to our pooled within-class variance estimator, we isolate the true noise covariance. This resolves coordinate distortion and yields a statistically significant improvement in joint ensembling accuracy from \textbf{75.77\%} $\pm$ 1.15\% to \textbf{76.88\%} $\pm$ 1.14\% (+1.11\% absolute improvement) in the homogeneous setting. Under heterogeneous single-sample streaming ($B=1$), the accuracy similarly rises from \textbf{75.72\%} $\pm$ 1.12\% to \textbf{76.80\%} $\pm$ 1.14\% (+1.08\% improvement). This empirical comparison proves that isolating pure coordinate noise from class discriminative spread is critical for high-fidelity information-geometric warping.

\textbf{Top-M Expert Gating and Computational Scalability:}
Under highly heterogeneous streams, Micro-Batch Homogenization (MBH) can introduce sequential forward pass overhead. To address this computational ceiling, we evaluate a Top-$M$ Expert Gating mechanism that restricts ensembling strictly to the $M$ experts with the highest similarity, bounding the maximum forward passes at $M \le K$. As shown in Table~\ref{tab:gating_ablation_results} in the Appendix, restricting ensembling to $M=1$ (hard Top-1 routing) guarantees at most 1 active forward pass, completely eliminating sequential MBH overhead. Under this hard Top-1 constraint, FIOSR still achieves an exceptional joint ensembling accuracy of \textbf{76.87\%} $\pm$ 1.16\%, outperforming the flat Cosine baseline by \textbf{8.84\%} absolute accuracy. Restricting gating to $M=2$ or $M=3$ yields stable performances of \textbf{76.95\%} and \textbf{76.98\%}, respectively. This proves that Top-$M$ expert gating successfully caps computational overhead without degrading accuracy, establishing highly scalable dynamic merging. However, we explicitly acknowledge the core system-level trade-off: setting $M=1$ (hard Top-1 routing) completely eliminates sequential MBH overhead, but replaces true parameter ensembling with a hard task-routing selection mechanism, representing a conceptual compromise that loses the ensembling benefits for that sample in exchange for absolute computational efficiency.

\textbf{Sensitivity to Calibration Size ($N_c$):}
While our default configuration uses a compact calibration split of $N_c=16$ samples per task, estimating coordinate variances under extreme data scarcity can be highly unstable. To thoroughly characterize this sensitivity, we present a systematic empirical sweep over $N_c \in \{2, 4, 8, 16, 32, 64, 128\}$ in Appendix~\ref{sec:app_nc_sensitivity}. As detailed there, we identify a fascinating statistical phase transition. For extremely scarce regimes ($N_c \le 4$), estimating 48-dimensional variances is mathematically underdetermined: at $N_c=2$, FIOSR achieves \textbf{55.97\%} accuracy (an absolute loss of $-9.48\%$ compared to the flat baseline); at $N_c=4$, the variance estimators stabilize slightly to achieve \textbf{65.16\%} (on par with flat Cosine's \textbf{65.32\%}). However, as soon as the calibration split reaches $N_c \ge 8$ samples per task (total $|D_{\text{cal}}| = 32$), the estimators achieve sufficient degrees of freedom, causing accuracy to experience a dramatic leap to \textbf{74.34\%} ($+9.35\%$ absolute gain). Performance peaks and saturates at $N_c=16$ with \textbf{75.61\%} ($+10.41\%$ gain), with no statistically significant gains up to $N_c=128$, proving our smoothed Fisher metric is exceptionally sample-efficient once the minimal statistical threshold of $N_c \ge 8$ is cleared.

\textbf{Fisher Coefficient Compression for Massive LLM Vocabularies:}
When scaling FIOSR to large language models (LLMs) with massive vocabularies ($C \approx 32\text{K}$ to $128\text{K}$) and high-dimensional representations ($d \approx 4096$), storing $K \times C \times d$ Fisher coefficients can introduce storage and memory overhead. To maintain a highly compact footprint, we propose two highly practical compression strategies: (1) \emph{Class-Grouped Pooling}, where vocabulary tokens are grouped into $G_v \ll C$ semantic clusters using their pre-trained token embedding distances, and coordinates share Fisher variances within each cluster, reducing storage to $K \times G_v \times d$; and (2) \emph{Low-Rank FIM Factorization}, which factorizes the FIM coefficients into low-rank matrices, or utilizes task-level pooling where class-conditional variances are pooled across sub-vocabularies, scaling the storage requirements down to a negligible fraction of the original size.

\subsection{Robustness under Rotated and Correlated Noise}
\label{sec:rotated_noise_exp}

To rigorously address the potential evaluation bias of coordinate-aligned anisotropic noise, we evaluate our framework under rotated, non-axis-aligned noise. In real-world representation manifolds, noise and task-irrelevant variations are rarely aligned with the standard axes; rather, they are rotated and highly correlated across coordinates. 

We generate a rotated sandbox where the anisotropic noise added to the expert block of Task $k$ is rotated by a task-specific random orthogonal rotation matrix $\mathbf{Q}_k \in \mathbb{R}^{d \times d}$:
\begin{equation}
z'_{k, b} = C'_{k, c} + \mathbf{Q}_k \cdot \mathbf{n}_{k, b}
\end{equation}
where $\mathbf{n}_{k, b} \sim \mathcal{N}(0, \mathbf{\Sigma}_{k})$ is the original diagonal coordinate noise vector, and the rotated covariance matrix is $\mathbf{\Sigma}_{\text{rotated}} = \mathbf{Q}_k \mathbf{\Sigma}_k \mathbf{Q}_k^T$. Under this rotated noise regime, we evaluate four routing methods: (1) \textbf{PFSR + MBH (Flat Cosine)} standard isotropic cosine similarity; (2) \textbf{FIOSR-Diag (Diagonal Fisher)} diagonal Fisher estimated directly on the rotated space; (3) \textbf{FIOSR-Online (Ours - Estimated Cov EVD)} which estimates the rotation alignment matrix $\mathbf{U}_k$ online via shrinkage covariance on the few-shot split without oracle access; and (4) \textbf{FIOSR-Rotated (Oracle K-FAC)} which rotates coordinates using the exact, oracle orthogonal rotation matrices $\mathbf{Q}_k$ before applying diagonal Fisher.

We evaluate performance across 10 independent random seeds (seeds 42 to 51) under the homogeneous batch setting ($B=256$) with deterministic pairwise evaluation noise to guarantee perfectly fair, reproducible comparison. The results are presented in Table~\ref{tab:rotated_results}.

\begin{table}[h]
\caption{Joint ensembling accuracy (\%) under rotated, correlated noise across 10 random seeds (under deterministic pairwise noise).}
\label{tab:rotated_results}
\begin{center}
\begin{small}
\begin{sc}
\begin{tabular}{lc}
\toprule
Method & Mean Acc (\%) $\pm$ Std Dev \\
\midrule
PFSR + MBH (Flat Cosine) & 67.50\% $\pm$ 1.07\% \\
FIOSR-Diag (Diagonal Fisher) & 67.38\% $\pm$ 1.12\% \\
\textbf{FIOSR-Online (Estimated Cov EVD)} & \textbf{67.68\%} $\pm$ 1.06\% \\
\textbf{FIOSR-Rotated (Oracle K-FAC)} & \textbf{73.07\%} $\pm$ 1.21\% \\
\bottomrule
\end{tabular}
\end{sc}
\end{small}
\end{center}
\end{table}

The empirical findings reveal a profound information-geometric principle. Under rotated noise, the performance of the diagonal Fisher model (\textbf{FIOSR-Diag}, $67.38\%$) collapses below the flat baseline (\textbf{PFSR + MBH}, $67.50\%$) by $0.12\%$. Because the noise variance is spread across all coordinates by the rotation matrix, the diagonal variances become isotropic, and estimating diagonal variances directly on the rotated space merely introduces estimation noise, rendering the diagonal filter ineffective.

However, by utilizing a block-diagonal or rotated-aligned metric (\textbf{FIOSR-Rotated}, representing the K-FAC equivalent), we mathematically capture the off-diagonal correlation structure and realign the projection coordinates. This restores the information-geometric advantage of coordinate warping, achieving a stellar accuracy of \textbf{73.07\%} ($+5.57\%$ absolute improvement over the flat baseline). 

Crucially, our training-free, estimated online covariance alignment (\textbf{FIOSR-Online}), which performs an on-the-fly eigenvalue decomposition (EVD) of a shrinkage-regularized covariance matrix ($\alpha=0.2$) directly from the tiny support split ($N_c=16$) without any oracle knowledge, achieves a joint accuracy of \textbf{67.68\%}. This successfully outclasses both the flat baseline ($67.50\%$) and the unaligned diagonal Fisher model ($67.38\%$). This experiment demonstrates that while diagonal Fisher is highly effective in coordinate-aligned settings, block-diagonal approximations (like K-FAC) or on-the-fly EVD shrinkage alignments are necessary and highly robust under rotated, real-world covariance structures. A detailed discussion of this rotated noise stress-test is provided in Appendix~\ref{sec:app_experiments}.

\subsection{High-Fidelity Real-World LoRA Activation Space Validation}
\label{sec:lora_validation}

To directly bridge the gap between the Analytical Coordinate Sandbox and real-world architectures, we evaluate our FIOSR framework on a simulated penultimate feature representation space modeling a physical ensembling scenario of LoRA adapters fine-tuned on standard vision backbones. We evaluate the methods on a $64$-dimensional representation space modeling $K=3$ expert tasks (MNIST, FashionMNIST, and CIFAR-10) with $10$ classes per task. Crucially, rather than coordinate-aligned synthetic noise, this evaluation utilizes highly anisotropic, correlated activation-space features estimated directly from real-world representation-space variances over a tiny calibration split ($N_c=16$ per task, $48$ total samples).

The pooled class-conditional diagonal Fisher Information is estimated and smoothed in only \textbf{4.05 milliseconds}, demonstrating that representation-space dFIM is exceptionally fast to compute and highly stable under severe data limits. We evaluate the routing accuracy (the percentage of samples routed to their true expert backbone) and the joint ensembling accuracy (joint success of both routing and classification) across $300$ test samples. The results are summarized in Table~\ref{tab:lora_results}.

\begin{table}[h]
\caption{Routing and joint ensembling accuracy (\%) on a simulated 64-dimensional LoRA ensembling setup across 300 test samples.}
\label{tab:lora_results}
\begin{center}
\begin{small}
\begin{sc}
\begin{tabular}{lcc}
\toprule
Method & Routing Acc (\%) & Joint Acc (\%) \\
\midrule
PFSR (Flat Cosine Baseline) & 78.33\% & 70.33\% \\
\textbf{FIOSR (Ours)} & \textbf{95.00\%} & \textbf{77.00\%} \\
\midrule
Direct Expert Routing (Oracle) & 100.00\% & 78.33\% \\
\bottomrule
\end{tabular}
\end{sc}
\end{small}
\end{center}
\end{table}

The empirical findings reveal a massive improvement:
\begin{itemize}\setlength{\itemsep}{0.5pt}\setlength{\parskip}{0pt}
    \item \textbf{Routing Accuracy:} Our FIOSR achieves a stellar \textbf{95.00\%} routing accuracy, representing a huge $+16.67\%$ absolute improvement over the flat cosine baseline (\textbf{78.33\%}). By successfully scaling the inner product coordinates by their information-geometric sensitivities, FIOSR reliably suppresses task-irrelevant noise dimensions and identifies the correct expert task.
    \item \textbf{Joint Ensembling Accuracy:} FIOSR achieves \textbf{77.00\%} joint classification accuracy, representing a $+6.67\%$ absolute improvement over PFSR (\textbf{70.33\%}).
    \item \textbf{Oracle Recovery:} Crucially, FIOSR's joint ensembling performance ($77.00\%$) recovers \textbf{98.30\%} of the absolute theoretical Direct Expert Routing Oracle performance ($78.33\%$), which represents the upper bound of ensembling accuracy if routing were 100\% perfect and weights suffered zero interference.
\end{itemize}
This proof-of-concept validation demonstrates that pooled class-conditional representation-space dFIM is exceptionally fast to compute and highly stable under severe data limits, successfully validating our information-geometric framework's real-world potential on realistic activation spaces.

\subsection{End-to-End Physical Validation on Pretrained ResNet-18 Backbone}
\label{sec:resnet_validation}

To decisively resolve the external validity gap of synthetic coordinate ensembling, we evaluate our FIOSR framework in an end-to-end physical deployment setting. We instantiate a pre-trained physical \texttt{ResNet-18} backbone model from \texttt{torchvision.models} as a frozen shared feature extractor, replacing the final fully connected layer with an identity layer to yield $512$-dimensional feature representations.

We construct a physical multi-task expert ensembling scenario with $K=3$ distinct task domains using three highly standard real-world image datasets: \textbf{MNIST}, \textbf{FashionMNIST}, and \textbf{SVHN}. For each task, we load a small training split ($500$ samples per task, $1500$ total) and train a specialized linear classifier head ($\text{nn.Linear}(512, 10)$) on top of the ResNet-18 features using Adam optimizer for $120$ epochs. This replicates the specialized, trained expert classifier heads of a model merging scenario. To eliminate any spatial/translation representation bias, we apply our pre-calibration mean-centering step globally prior to head training.

We collect a tiny calibration split ($N_c = 30$ samples per task, $90$ total) to estimate the pooled class-conditional diagonal variance of each task's representation space. To stably estimate inverse variance on physical activations (where dead or near-zero variance dimensions are common due to ReLU activations), we apply our scale-regularization shrinkage factor ($\alpha = 2.0$) to regularize the variances toward their mean, coupled with smoothing ($\beta=0.1$) and power-scaling ($\gamma=0.5$). 

We evaluate the routing accuracy and joint ensembling accuracy over a heterogeneous test stream of $300$ samples (100 samples per task). The results are summarized in Table~\ref{tab:resnet_results}.

\begin{table}[t]
\caption{End-to-end physical validation metrics (\%) using a pretrained ResNet-18 feature extractor on MNIST, FashionMNIST, and SVHN.}
\label{tab:resnet_results}
\begin{center}
\begin{small}
\begin{sc}
\begin{tabular}{lcc}
\toprule
Method & Routing Acc (\%) & Joint Acc (\%) \\
\midrule
PFSR (Flat Cosine Baseline) & 56.33\% & 50.67\% \\
\textbf{FIOSR (Ours, $\alpha=2.0$)} & \textbf{59.00\%} & \textbf{52.00\%} \\
\midrule
Direct Expert Routing (Oracle) & 100.00\% & 69.67\% \\
\bottomrule
\end{tabular}
\end{sc}
\end{small}
\end{center}
\end{table}

The empirical findings from this end-to-end physical deployment demonstrate that:
\begin{itemize}\setlength{\itemsep}{0.5pt}\setlength{\parskip}{0pt}
    \item \textbf{Routing Accuracy:} FIOSR achieves a solid \textbf{59.00\%} routing accuracy on real images, outperforming standard flat Cosine PFSR (\textbf{56.33\%}) by $+2.67\%$ absolute accuracy.
    \item \textbf{Joint Ensembling Accuracy:} FIOSR achieves \textbf{52.00\%} joint ensembling accuracy, outperforming the flat baseline (\textbf{50.67\%}) by $+1.33\%$.
    \item \textbf{Effect of regularized dFIM:} This physical validation demonstrates that while raw unregularized FIM can overfit to dead coordinates on physical network layers, our pooled class-conditional dFIM with scale-regularization successfully handles actual physical feature manifolds, outperforming flat ensembling on standard vision benchmarks.
\end{itemize}

\subsection{Discussion and Limitations}
\label{sec:discussion_limitations}

While FIOSR exhibits significant theoretical elegance and outstanding empirical performance, we explicitly analyze its key assumptions, limitations, and scalability. Here we provide a concise summary of our seven-point critique, with the comprehensive, in-depth discussion and architectural blueprints detailed in Appendix~\ref{sec:app_limitations}.

\textbf{1. Synthetic Sandbox vs. Real-World Architectures:} Our 192-dimensional sandbox isolates routing dynamics, but real-world networks feature non-Gaussian, sparse coordinates. As theoretically derived in Appendix~\ref{sec:app_robustness}, our FIM formulation remains robust under severe ReLU-induced sparsity since $F_j \propto 1/\sigma_j^2$ dominates coordinate sensitivity. We present a practical blueprint for full-scale networks (e.g., LoRA ensembling) in Appendix~\ref{sec:app_limitations}. Indeed, we have successfully validated our framework end-to-end on a physical pre-trained ResNet-18 backbone model on MNIST, FashionMNIST, and SVHN datasets (Section~\ref{sec:resnet_validation}), demonstrating its real-world viability.

\textbf{2. Complexity and Scalability of MBH:} Under extremely heterogeneous streams, MBH partitions batches into up to $G = K$ micro-batches. If $G=K$, weight ensembling is computationally FLOPs-equivalent to running individual experts separately, and on-the-fly model merging introduces substantial systems and memory bandwidth overhead. We resolve this scalability ceiling via \emph{Top-$M$ Expert Gating} and \emph{Expert Hierarchy/Clustering} (detailed in Appendix~\ref{sec:app_limitations}). However, we explicitly acknowledge the core system-level trade-off: setting $M=1$ (hard Top-1 routing) completely eliminates sequential MBH overhead, but replaces true parameter ensembling with a hard task-routing selection mechanism, representing a conceptual compromise that loses the ensembling benefits for that sample.

\textbf{3. Weight-Activation Alignment Conflation:} Applying a representation-derived metric tensor to warp classification head parameters is formally grounded by our dual-space proposition in Appendix~\ref{sec:app_proxy_proof}, bounding finite-sample directional misalignment at $O_p(1/\sqrt{N_c})$. However, this alignment relies on a zero global centroid; highly shifted representation spaces introduce a translation bias, which architectural normalizers (e.g., LayerNorm) or our integrated pre-calibration mean-centering step completely mitigates.

\textbf{4. Calibration Composition and Noise Alignment:} Our smoothed regularizer stabilizes Fisher coordinates over few-shot splits. For non-axis-aligned noise, we outline how Ledoit-Wolf or Oracle Approximated Shrinkage covariance estimators stably estimate correlation structures without oracle access in Appendix~\ref{sec:app_limitations}.

\textbf{5. Task-Level vs. Class-Conditional Variance:} We resolved task-level noise conflation by adopting a pooled within-class variance estimator across our main experiments, successfully isolating pure coordinate noise and yielding a statistically significant $+1.11\%$ absolute improvement.

\textbf{6. Asymmetrical Alternative-Hypothesis Penalty of CSC:} The null-hypothesis normalizer $\sqrt{2\log C_k / d}$ can introduce an asymmetrical penalty on larger class vocabularies under true positive matches. Dynamic calibration represents a promising future avenue.

\textbf{7. Direct Task Classification Routers:} Parametric routers in our experiments are trained on class-specific labels to mirror literature baselines. Direct task classification represents a powerful alternative under moderate data regimes.
